\section{\ours: Training-Free and Adaptive Guard for Safe Text-to-Image and Video Generation}

% \hy{I suggest to ref the figure 2 here, and also leverage figure 2 when introduing the method. }
% \updated{\textbf{Limitations in Existing Works.}} 
% Recent approaches~\citep{gandikota2024unified,gong2024reliable,lu2024mace,chavhan2024conceptprune,yang2024pruning}  have demonstrated the effectiveness of weight modification through unlearning or model editing to prevent the generation of harmful (e.g., pornography, self-harm, violence), biased (e.g., racial or social stereotypes, ageism), or otherwise undesirable (e.g., public, copyrighted) visual content in T2I models. 
% However, these methods have limited flexibility because they 1) require storing individual model weights for each concept to be removed, 2) reduce the backbone model's generative capabilities through unlearning, and 3) necessitate distinct solutions for safe generation across different models (i.e., modified model weights).

% \updated{\textbf{\ours{}.}} To address this, w
We propose \ours{}, a training-free, yet adaptive remedy for safe T2I and T2V generation (\Cref{fig:concept}). 
We first determine the trigger tokens that can potentially induce toxicity based on the proximity between masked input prompt embeddings and toxic concept subspace (\Cref{sec:subsec:concept_subspace}). 
{% Detected embeddings are projected orthogonally to the concept subspace while staying in the input space
The detected trigger token embeddings are projected onto a subspace orthogonal to the toxic concept subspace, while maintaining them within the input space
(\Cref{sec:subsec:orthogonal_modification}).
\ours{} automatically adjusts the number of denoising timesteps conditioned on the text inputs through a self-validating filtering mechanism (\Cref{sec:subsec:selfvalid}). }Additionally, 
we introduce an adaptive re-attention strategy in the latent space during the de-noising process, facilitating a robust joint filtering mechanism across both text and visual embedding spaces (\Cref{sec:subsec:latent_attention}).
Finally, we extend our approach to high-resolution models like SDXL~\citep{podell2023sdxl}, DiT-based image diffusion (SD-v3)~\citep{peebles2023scalable}, and representative text-to-video generative models such as ZeroScopeT2v~\citep{2023zeroscope} and CogVideoX~\citep{yang2024cogvideox} (\Cref{sec:subsec:video_generation}).

\begin{figure}[t]
    \centering
    \vspace{-0.2in}\includegraphics[width=0.9\linewidth]{assets/concept-figure.pdf}
\vspace{-0.1in}
\caption{\textbf{\ours{} framework}. Based on proximity analysis between the masked token embeddings and the toxic subspace~$\mathcal{C}$, we detect unsafe tokens and project them into orthogonal to the toxic concept (in \cpink{red}), but still be in the input space~$\mathcal{I}$ (in \cgreen{green}). SAFREE adaptively controls the filtering strength in an input-dependent manner, which also regulates a latent-level re-attention mechanism. Note that our approach can be broadly applied to various image and video diffusion backbones.}
    \label{fig:concept}\vspace{-0.1in}
\end{figure}


\subsection{Adaptive Token Selection based on Toxic Concept Subspace Proximity}\label{sec:subsec:concept_subspace}


% \textbf{Classifier-free guidance and negative prompts.} 
Random noise $\epsilon_0\sim\mathcal{N}(0,\bsy{I})$ sampled from a Gaussian distribution can lead to the generation of unsafe or undesirable images in diffusion models, primarily due to inappropriate semantics embedded in the text, which conditions the iterative denoising process~\citep{rombach2022high,ho2022classifier}. To mitigate this risk, recent studies~\citep{miyake2023negative,schramowski2023safe,ban2024understanding} have demonstrated the effectiveness of using negative prompts. In this approach, the model aims to predict the refined noise from $\epsilon_0$ over several autoregressive denoising steps, synthesizing an image conditioned on the input (i.e., the input text prompt).
The denoising process of diffusion models, parameterized by $\bsy{\theta}$, at timestep $t$ follows the classifier-free guidance approach:
\begin{equation}
\epsilon_t=\left(1+\omega\right)\epsilon_{\bsy\theta}\left(\bsy{z}_t,\bsy{p}\right)-\omega\epsilon_{\bsy\theta}\left(\bsy{z}_t,\emptyset\right),
\end{equation}
where $\omega$ is a hyperparameter controlling the guidance scale. 
$\bsy{p}$ and $\emptyset$ denote the embedding of the input prompt and null text, respectively. 
The negative prompt is applied by replacing $\emptyset$ with the embedding of the negative prompt. We note that even if the input prompts are adversarial and unreadable to humans, such as those generated by adversarial attack methods through prompt optimization~\citep{tsai2024ring,yang2024mma,chin2024prompting4debugging}, they are still encoded within the same text embedding space, like CLIP~\citep{radford2021learning}, highlighting the importance and necessity of the feature embedding level unsafe concepts identification for safeguard.
To this end, we propose to detect token embeddings that trigger inappropriate image generation and transform them to be distant from the toxic concept subspace: $\mathcal{C}=\left[\bsy{c}_0; \bsy{c}_1; ...; \bsy{c}_{K-1}\right]\in\mathbb{R}^{D\times K}$, which represents the embedding matrix that denotes the toxic subspace, where each column vector \updated{$\bsy{c}_k$} corresponds to the embedding of the relevant text associated with the $k$-th user-defined toxic concept, such as \textit{Sexual Acts} or \textit{Pornography} for \textit{Nudity} concept.
% a text embedding extracted derived from a keyword from a set of keywords associated with toxic concepts.

\updated{\textbf{Detecting Trigger Tokens Driving Toxic Outputs.}} To assess the relevance of specific tokens in the input prompt to the toxic concept subspace, we design a pooled input embedding \(\overline{\bsy{p}}_{\backslash i} \in \mathbb{R}^{D}\) that averages the token embeddings in \(\bsy{p}\) while masking out the \(i\)-th token.
Suppose $\bsy{z}$\updated{$\in\mathbb{R}^{K}$}  
\updated{be the vector of coefficients (i.e., the projection coordinates)}
% be a mapped embedding 
of $\overline{\bsy{p}}_{\backslash i}$ onto $\mathcal{C}$, it satisfies, 
\begin{equation}
\begin{split}
{\mathcal{C}}^\top\left(\overline{\bsy{p}}_{\backslash i}-\bsy{z}\mathcal{C}\right)=0,\quad\quad\quad\quad
\bsy{z}=\left({\mathcal{C}}^\top\mathcal{C}\right)^{-1}{\mathcal{C}}^\top \overline{\bsy{p}}_{\backslash i}.
\end{split}
\end{equation}
We estimate the conceptual proximity of a token in the input prompt with $\mathcal{C}$ by computing the distance between the pooled text embedding obtained after masking out (i.e., removing) the corresponding token and $\mathcal{C}$. 
The residual vector \(\bsy{d}_{\backslash i}\), which is the component of \(\overline{\bsy{p}}_{\backslash i}\) orthogonal to the subspace \(\mathcal{C}\) is then formulated as follows:
\begin{equation}
\begin{split}\label{eq:residual}
\bsy{d}_{\backslash i}=\overline{\bsy{p}}_{\backslash i}-\mathcal{C}\bsy{z}
&=\left(\bsy{I}-\mathcal{C}\left({\mathcal{C}}^\top\mathcal{C}\right)^{-1}{\mathcal{C}}^\top\right)\overline{\bsy{p}}_{\backslash i}\\
&=\left(\bsy{I}-\bsy{P}_{\mathcal{C}}\right) \overline{\bsy{p}}_{\backslash i},~~~~\text{where}~~\bsy{P}_{\mathcal{C}}=\mathcal{C}\left({\mathcal{C}}^\top\mathcal{C}\right)^{-1}{\mathcal{C}}^\top
\end{split}
\end{equation}
and $\bsy{I}\in\mathbb{R}^{D \times D}$ denotes the identity matrix (See~\Cref{fig:concept} middle left).
A longer residual vector distance indicates that the removed token in the prompt is more strongly associated with the concept we aim to eliminate. In the end, we derive a masked vector $\bsy{m}\in\mathbb{R}^{N}$ (where $N$ denotes the token length) to identify tokens related to the target concept, allowing us to subtly project them within the input token subspace while keeping them distant from the toxic concept subspace.
We obtain a set of distances of masked token embeddings $\overline{\bsy{p}}_{\backslash i}, i\in[0,N-1]$ to the concept subspace, $D(\bsy{p}|\mathcal{C})$, and select tokens for masking by evaluating the disparity between each token's distance and the average distance of the set, excluding the token itself:
\begin{equation}
\begin{split}
D(\bsy{p}|\mathcal{C})&=\left[\|\bsy{d}_{\backslash 0}\|_2, \|\bsy{d}_{\backslash 1}\|_2, ..., \|\bsy{d}_{\backslash N-1}\|_2\right],\\
m_i&=\begin{cases} 1 & \text{if}~~\|\bsy{d}_{\backslash i}\|_2 > \left(1+\alpha\right)\cdot\text{mean}\left(D(\bsy{p}|\mathcal{C}).\text{\updated{delete}}\left(i\right)\right),\\
                     0 &  \text{otherwise},
       \end{cases}
\end{split}
\end{equation}
where $\alpha$ is a non-negative hyperparameter that controls the sensitivity of detecting concept-relevant tokens. $X.\text{delete}(i)$ denotes an operation that produces a list $X$ removing the $i$-th item. We set $\alpha=0.01$ for all experiments in this paper, demonstrating the robustness of our approach to $\alpha$ across T2I generation tasks with varying concepts. We project the detected token embedding (i.e., $m_i=1$) to safer embedding space (See \Cref{sec:subsec:orthogonal_modification}).

\subsection{Safe Generation via Concept Orthogonal Token Projection}\label{sec:subsec:orthogonal_modification}

We aim to project toxic concept tokens into a safer space to encourage the model to generate appropriate images. However, directly removing or replacing these tokens with irrelevant ones, such as random tokens or replacing the token embeddings with null embeddings, disrupts the coherence between words and sentences, compromising the quality of the generated image to the safe input prompt, particularly when the prompt is unrelated to the toxic concepts. To address this, we propose projecting the detected token embeddings into a space orthogonal to the toxic concept subspace while keeping them within the input space to ensure that the integrity of the original prompt is preserved as much as possible.
We begin by formalizing the input space \(\mathcal{I}\) using pooled embeddings from masked prompts as described in~\Cref{sec:subsec:concept_subspace}, such that
$\mathcal{I}=\left[\overline{\bsy{p}}_{\backslash 0}; \overline{\bsy{p}}_{\backslash 1}; ...; \overline{\bsy{p}}_{\backslash N-1}\right]\in\mathbb{R}^{D\times N}$. 
Given the projection matrix into input space $\mathcal{I}$ formulated by $\bsy{P}_\mathcal{I}=\mathcal{I}\left({\mathcal{I}}^\top\mathcal{I}\right)^{-1}{\mathcal{I}}^\top$ (derived by \Cref{eq:residual}), we perform selective detoxification of input token embeddings based on the obtained token masks that project assigned tokens into $\bsy{P}_{\mathcal{I}}$ and to be orthogonal to $\bsy{P}_{\mathcal{C}}$:
\begin{equation}
\begin{split}
\bsy{p}_{proj}&=\bsy{P}_{\mathcal{I}}\left(\bsy{I}-\bsy{P}_{\mathcal{C}}\right)\bsy{p},\\
\bsy{p}_{safe}&=\bsy{m}\odot\bsy{p}_{proj}+(\bsy{1}-\bsy{m})\odot\bsy{p},
\end{split}
\end{equation}
where $\odot$ indicates an element-wise multiplication operator. \updated{That is, for the $i$-th token, we use the projected safe embeddings $\bsy{p}_{proj, i}$ only if it is detected as a toxic token ($m_i\odot\bsy{p}_{proj, i}, m_i=1$); otherwise, we retain the original (safe) token embeddings $\bsy{p}_i$, since $(\bsy{1}-\bsy{m_i})\odot\bsy{p}_i, m_i=0$.}


\subsection{\updated{Adaptive Control of Safeguard Strengths with Self-validating Filtering}}\label{sec:subsec:selfvalid}
While our approach so far adaptively controls the number of token embeddings to be updated, it sometimes lacks flexibility in preserving the original generation capabilities for content outside the target concept. Recent observations~\citep{kim2024race,ban2024understanding} suggest that different denoising timesteps in T2I models contribute unevenly to generating toxic or undesirable content. Based on this insight, we propose a self-validating filtering mechanism during the denoising steps of the diffusion model that automatically adjusts the number of denoising timesteps conditioned on the obtained embedding (middle in \Cref{fig:concept}). This mechanism amplifies the model's filtering capability when the input prompt is deemed undesirable, while approximating the original backbone model's generation for safe content. In the end, our updated input text embedding $\bsy{p}_{safree}$ at a different denoising step $t$ is determined as follows:
\begin{equation}
\begin{split}\label{eq:self-valid}
t'=\gamma\cdot\text{sigmoid}(1-\text{cos}(\bsy{p}, \bsy{p}_{proj})),\quad\quad
\bsy{p}_{safree}=\begin{cases} \bsy{p}_{safe} & \text{if}~~t\leq\text{round}(t'),\\
                     \bsy{p} &  \text{otherwise},
       \end{cases}
\end{split}
\end{equation}
where $\gamma$ is hyperparameter ($\gamma=10$ throughout the paper) and \textit{cos} represents cosine similarity. $t'$ denotes the self-validating threshold to determine the number of denoising steps applying to the proposed safeguard approach. Specifically, we adopt the cosine distance between the original input embedding $\bsy{p}$ and the projected embedding $\bsy{p}_{proj}$ to compute $t'$.  A higher similarity indicates that the input prompt has been effectively disentangled from the toxic target concept to be removed. 

\subsection{Adaptive Latent Re-attention in Fourier Domain}\label{sec:subsec:latent_attention}
Recent literature~\citep{mao2023guided,qi2024not,ban2024crystal,sun2024spatial} has demonstrated that the initial noise sampled from a Gaussian distribution significantly impacts the fidelity of T2I generation in diffusion models. To further guide these models in creating content while suppressing the appearance of inappropriate or target concept semantics, we propose a novel visual latent filtering strategy during the denoising process. \cite{si2024freeu} note that current T2I models frequently experience oversmoothing of textures during the denoising process, resulting in distortions in the generated images.
Building on this insight, we suggest an adaptive re-weighting strategy using spectral transformation in the Fourier domain. At each timestep, we initially perform a Fourier transform on the latent features, conditioned on the initial prompt $\bsy{p}$ (which may incorporate unsafe guidance) and our filtered prompt embedding $\bsy{p}_{safree}$.
The low-frequency components typically capture the global structure and attributes of an image, encompassing its overall context, style, and color.
In this context, we reduce the influence of low-frequency features, which are accentuated by our filtered prompt embedding, while preserving the visual regions that are more closely aligned with the original prompt to avoid excessive oversmoothing. Let $h(\cdot)$ be a latent feature, to achieve this, we attenuate the low-frequency features in $h(\bsy{p}_{safree})$ by a scalar $s$ when their values are lower in magnitude than those from $\bsy{p}$:
\begin{equation}
\begin{split}\label{eq:fft}
\mathcal{F}({\bsy{p}})&=\bsy{b}\odot\text{FFT}(h(\bsy{p})),~~~\mathcal{F}({\bsy{p}_{safree}})=\bsy{b}\odot\text{FFT}(h(\bsy{p}_{safree})),
\end{split}
\end{equation}
\begin{equation}
\begin{split}\label{eq:fft-select}
\mathcal{F}'_i&=\begin{cases} s \cdot \mathcal{F}({\bsy{p}_{safree}})_i & \text{if}~~ \mathcal{F}({\bsy{p}_{safree}})_i > \mathcal{F}({\bsy{p}})_i,\\
                     \mathcal{F}({\bsy{p}_{safree}})_i &  \text{otherwise}.
       \end{cases}
\end{split}
\end{equation}
where $s < 1$, $\bsy{b}$ represents the binary masks corresponding to the low-frequency components (i.e., the middle in the width and height dimension), and FFT denotes a Fast Fourier Transform operation. We first obtain low frequency features from $h(\bsy{p})$ and $h(\bsy{p}_{safree})$ in~\Cref{eq:fft}. This reduces the oversmoothing effect in safe visual components, encouraging the generation of safe outputs without emphasizing inappropriate content. This process is enabled by obtaining the refined features $h'$ via inverse FFT, $h'=\text{IFFT}(\mathcal{F}')$, as described in~\Cref{eq:fft-select}. Note that this equation doesn't affect the original feature if $t > \text{round}(t')$ since $\mathcal{F}({\bsy{p}_{safree}})_{i}==\mathcal{F}({\bsy{p}})_i$, allowing automatic control of filtering capability through self-validated filtering.



\subsection{\ours{} for Advanced T2I Models and Text-to-Video Generation}\label{sec:subsec:video_generation}
Unlike existing unlearning-based methods limited to specific models or tasks, \ours{} is architecture agnostic and can be integrated across diverse backbone models without model modifications, offering superior versatility in safe generation. 
This flexibility is enabled by \textit{concept-orthogonal, selective token projection} and \textit{self-validating adaptive filtering}, allowing \ours{} to work across a wide range of generative models and tasks. 
It operates seamlessly with models beyond SD v-1.4, like, SDXL~\citep{podell2023sdxl} and SD-v3~\citep{2024sdv3} in a zero-shot, training-free manner, and extends its applicability to text-to-video (T2V) generation models like ZeroScopeT2V~\citep{2023zeroscope} and CogVideoX~\citep{yang2024cogvideox}, making it highly flexible as a series of plug-and-play modules. 
We present both qualitative and quantitative results showing its effectiveness across various model backbones (UNet~\citep{ronneberger2015u} and DiT~\citep{peebles2023scalable}) and tasks (T2I and T2V generation) in~\Cref{sec:t2v}.
